%% toc: P15 157 157.7
% Source variants of the measure-theoretic foundations in Durrett, Probability: Theory and Examples, Version 5 (11 January 2019): pp. 2--3, 6--17, 53--54, 210, 214--215, 457--458.
\chapter{Weights, Generators, Extended Limits and Coding Variants}
\label{sec:p15-w02c-p15-002-source-variants}

The foundations in
\cref{sec:p15-measurable-probability,sec:p15-countable-product-measures,sec:p15-regular-conditional-laws}
follow Rick Durrett's \emph{Probability: Theory and Examples}, Version~5,
dated 11 January 2019. At several places that source states a variant of a
definition or theorem, a piece of notation, a terminology warning, a named
example, or an attribution that the earlier chapters replace by a
neighbouring form. This chapter states each of those source forms exactly
and proves how it relates to the form used elsewhere in the handbook.

\emph{Purpose.} To make every one of these source variants available in
its own terms, with the bridge to the handbook's form proved, so that a
reader working from the source can move between the two without a gap.

\emph{Inputs.} Countable atomic models (\cref{thm:p15-countable-atomic});
Borel sets and their generators (\cref{def:p15-borel-sets},
\cref{thm:p15-borel-generators}); measurable maps and their test by
generators (\cref{def:p15-measurable-map},
\cref{thm:p15-measurability-criterion}); extended real random variables
and their limits (\cref{def:p15-extended-random-variables},
\cref{thm:p15-measurable-limits}); the pi-lambda theorem
(\cref{thm:p15-measure-uniqueness}); Borel real codings
(\cref{def:p15-measurable-real-coding},
\cref{thm:p15-metric-borel-representation}); conditional expectation
(\cref{def:p15-general-conditional-expectation}) and regular conditional
laws (\cref{def:p15-regular-conditional-law}).

\emph{Conclusions.} The astragalus model with four unequal weights; the
open rational rays and the real-endpoint open rectangles as generators;
Durrett's conditions for a function on \(\mathbb R^d\) to define a
measure, his counterexample with a negative point mass, and his
Theorem~1.1.11 with the tail condition it needs; sampling a law by
applying the generalized inverse \(F^{-1}\) to a uniform variable; the notation
\(P\circ X^{-1}\) for the distribution; the term \emph{equal in
distribution} and the notation \(X\stackrel{d}{=}Y\); a discrete law whose
distribution function jumps at every rational number; the membership notation \(X\in\mathcal F\) for measurability and the
indicator--characteristic-function terminology; the extended line
\((\mathbf R^*,\mathcal R^*)\) with its order topology, random variables
defined on an event \(D\), the measurable convergence set \(\Omega_o\),
almost sure convergence without a prescribed limit, and the convention
\(X_\infty=\limsup_n X_n\); uniqueness of measures that agree on a
pi-system with an exhausting sequence of finite measure; the step in the
source's existence proofs that uses the image of a \emph{nice} map;
Borel's paradox in the source's framing; and the
source's account of why regular conditional distributions can fail to
exist; and the exact status of the source's existence theorems for
regular conditional distributions on nice spaces.

\emph{Relationships.} Each section names the earlier result it extends.
Every theorem below has its proof in the annex.

\section{Discrete weights: the astragalus}
\label{sec:p15-w02c-p15-002-astragali}

Durrett's discrete probability spaces assign a weight \(p(\omega)\geq0\)
to each point of a countable set, with total weight one, and sum weights
over an event \cite[Example~1.1.2, p.~2]{DurrettPTEAuthor2019}. In the
handbook this is the countable atomic model of
\cref{thm:p15-countable-atomic}. The source then gives a concrete case in
which equal weights would be unjustified
\cite[p.~3]{DurrettPTEAuthor2019}.

\begin{hbexample}{def:p15-w02c-p15-002-astragali}{The astragalus model}
Durrett describes the astragali as dice used in ancient Egypt, made from
the ankle bones of sheep \cite[p.~3]{DurrettPTEAuthor2019}. According to
his description a bone comes to rest in one of four ways: on the top side
(scored four points), on the bottom side (three points), on a flat narrow
part of the side (six points), or on the remaining, slightly rounded part
of the side (one point). He observes that nothing suggests these four
outcomes are equally likely, so the model needs four separate
probabilities \(p_1,p_3,p_4,p_6\).

\emph{Model.} Put \(\Omega=\{1,3,4,6\}\), the possible scores, with
\(\mathcal F=2^\Omega\), and choose weights
\[
 p_1,p_3,p_4,p_6\geq0,\qquad p_1+p_3+p_4+p_6=1 .
\]
By \cref{thm:p15-countable-atomic} the assignment
\(\mathbb P(A)=\sum_{i\in A}p_i\) is a probability measure on
\((\Omega,2^\Omega)\), and every probability measure on this space
arises from exactly one such weight vector, namely
\(p_i=\mathbb P(\{i\})\). The family of possible models is therefore the
set of all such vectors, a three-dimensional simplex; the uniform choice
\(p_i=1/4\) is one member of it and has no special standing.

\emph{Computation.} The event that the score is even is \(\{4,6\}\), with
probability \(p_4+p_6\). The event that the bone lands on its top or
bottom side is \(\{3,4\}\), with probability \(p_3+p_4\). For two
independent throws, the product model on \(\Omega^2\) with weights
\(p_ip_j\) (again a countable atomic model) gives the total score 7 the
probability \(2(p_1p_6+p_3p_4)\), since \(7=1+6=6+1=3+4=4+3\) and no
other ordered pair of scores sums to 7.

\emph{Status of the historical claim.} The description of the bones and
of their scoring is Durrett's report; the source cites no archaeological
evidence for it and gives no measured frequencies. The weights are
therefore free parameters of the model, not estimates. Nothing in the
model depends on the historical description beyond the number of
distinguishable outcomes.
\index{astragalus}\index{probability space!discrete}
\end{hbexample}

\section{Further generators of the Borel sets}
\label{sec:p15-w02c-p15-002-generators}

\Cref{thm:p15-borel-generators} gives closed rational rays, closed real
rays and rational open rectangles as generators. The source lists two
further choices for the test of \cref{thm:p15-measurability-criterion}:
the open rays \((-\infty,x)\) with rational \(x\) in one dimension, and
in dimension \(d\) the open rectangles with arbitrary real endpoints, or
the larger class of all open sets
\cite[Examples~1.3.2--1.3.3, p.~15]{DurrettPTEAuthor2019}.

\begin{hbtheorem}{thm:p15-w02c-p15-002-generator-variants}{proved=annex:p15-w02c-p15-002-generator-variants}{Open rational rays and real open rectangles generate the Borel sets}
\begin{enumerate}
\item \(\mathcal B(\mathbb R)=\sigma(\{(-\infty,x):x\in\mathbb Q\})\).
\item For every \(d\geq1\),
      \[
       \mathcal B(\mathbb R^d)
       =\sigma\bigl(\{(a_1,b_1)\times\cdots\times(a_d,b_d):
              -\infty<a_i<b_i<\infty\}\bigr),
      \]
      and \(\mathcal B(\mathbb R^d)\) is also generated by the class of
      all open sets, which contains these rectangles.
\item Consequently a function \(X:\Omega\to\mathbb R\) is a real random
      variable exactly when \(\{X<x\}\in\mathcal F\) for every rational
      \(x\), and a map \(X:\Omega\to\mathbb R^d\) is a random vector
      exactly when \(X^{-1}(R)\in\mathcal F\) for every open rectangle
      \(R\) with real endpoints.
\end{enumerate}
\end{hbtheorem}

The rational choices give countable generators; the real-endpoint
rectangles are uncountably many but are often the convenient ones to
check, for example when a set is described by strict inequalities. Each
class may be used in the test of \cref{thm:p15-measurability-criterion},
because each generates the target sigma-algebra.

\section{Measures on \texorpdfstring{\(\mathbb R^d\)}{R d} from distribution functions}
\label{sec:p15-w02c-p15-002-rd-measures}

In one dimension, \cref{thm:p15-stieltjes-construction} builds a Borel
measure from a nondecreasing right-continuous function, and
\cref{thm:p15-cdf-existence} realizes every admissible cumulative function
as the law of a random variable; these are the handbook's forms of the
source's Theorems~1.1.4 and~1.2.2, with the extension theorem of
\cref{thm:p15-semialgebra-extension,thm:p15-premeasure-extension}
\cite[Theorems~1.1.4, 1.2.2, pp.~3, 11]{DurrettPTEAuthor2019}. Durrett
then asks which functions on \(\mathbb R^d\) define measures, under the
heading \emph{Measures on \(\mathbf R^d\)}
\cite[pp.~6--9]{DurrettPTEAuthor2019}. This section states his
conditions, his counterexample and his Theorem~1.1.11, shows that the
theorem needs a stronger form of his condition~(iii), and proves the
corrected theorem.

\begin{hbdefinition}{def:p15-w02c-p15-002-rectangle-difference}{Finite rectangles, vertices and the difference \(\Delta_AF\)}
For \(x,y\in\mathbb R^d\) write \(x\leq y\) when \(x_i\leq y_i\) for all
\(i\), and \(y\downarrow x\) when each \(y_i\) decreases to \(x_i\). For
real \(a_i<b_i\), \(i=1,\ldots,d\), Durrett sets
\[
 A=(a_1,b_1]\times\cdots\times(a_d,b_d],\qquad
 V=\{a_1,b_1\}\times\cdots\times\{a_d,b_d\},
\]
calls \(A\) a \emph{finite rectangle} (no endpoint is infinite) and \(V\)
its set of vertices, and for \(v\in V\) and \(F:\mathbb R^d\to\mathbb R\)
defines
\[
 \operatorname{sgn}(v)=(-1)^{\#\text{ of }a\text{'s in }v},\qquad
 \Delta_AF=\sum_{v\in V}\operatorname{sgn}(v)F(v)
\]
\cite[p.~7]{DurrettPTEAuthor2019}. Thus \(\Delta_AF\) adds the values of
\(F\) at the vertices with an even number of lower endpoints and subtracts
those with an odd number. For \(d=1\), \(\Delta_{(a,b]}F=F(b)-F(a)\); for
\(d=2\),
\(\Delta_AF=F(b_1,b_2)-F(a_1,b_2)-F(b_1,a_2)+F(a_1,a_2)\).

Durrett's conditions on \(F:\mathbb R^d\to[0,1]\) are
\cite[pp.~6--7]{DurrettPTEAuthor2019}:
\begin{enumerate}
\item[(i)] \(F\) is nondecreasing: \(x\leq y\) implies \(F(x)\leq F(y)\);
\item[(ii)] \(F\) is right continuous:
      \(\lim_{y\downarrow x}F(y)=F(x)\);
\item[(iii)] if every coordinate of \(x_n\) decreases to \(-\infty\), then
      \(F(x_n)\downarrow0\); if every coordinate of \(x_n\) increases to
      \(\infty\), then \(F(x_n)\uparrow1\);
\item[(iv)] \(\Delta_AF\geq0\) for all finite rectangles \(A\).
\end{enumerate}
In the second clause of (iii) the printed text has \(x_n\uparrow-\infty\);
the parallel with the first clause and with the one-dimensional limit
\(F(x)\to1\) as \(x\to\infty\) shows that \(x_n\uparrow\infty\) is meant,
and that reading is used here.

The \emph{strong tail condition} (iii\(^*\)) is: for every coordinate
\(i\) and every \(x\in\mathbb R^d\),
\[
 F(x_1,\ldots,x_{i-1},t,x_{i+1},\ldots,x_d)\to0\quad(t\to-\infty),
\]
and \(F(x_n)\to1\) whenever every coordinate of \(x_n\) tends to
\(\infty\). Under (i), (iii\(^*\)) implies (iii).
\index{finite rectangle}\index{distribution function!on \(\mathbb R^d\)}
\end{hbdefinition}

\begin{hblemma}{thm:p15-w02c-p15-002-rectangle-inclusion-exclusion}{proved=annex:p15-w02c-p15-002-rectangle-inclusion-exclusion}{The difference \(\Delta_AF\) is the mass of the rectangle}
Let \(\mu\) be a measure on \((\mathbb R^d,\mathcal B(\mathbb R^d))\) with
\(F(x)\defeq\mu((-\infty,x_1]\times\cdots\times(-\infty,x_d])<\infty\) for
every \(x\). Then \(\mu(A)=\Delta_AF\) for every finite rectangle \(A\),
and \(F\) satisfies (i), (ii) and (iv). If \(\mu\) is a probability
measure, \(F\) also satisfies (iii\(^*\)), and hence (iii).
\end{hblemma}

For \(d=2\) the lemma is the identity Durrett writes out by inclusion and
exclusion \cite[p.~7]{DurrettPTEAuthor2019}. It shows that condition~(iv)
is necessary, and the next example, which is Durrett's, shows that
(i)--(iii) do not imply it.

\begin{hbexample}{def:p15-w02c-p15-002-negative-mass}{Durrett's counterexample: a negative mass at \((1,1)\)}
Let \(d=2\) and
\[
 F(x_1,x_2)=
 \begin{cases}
  1,&x_1,x_2\geq1,\\
  2/3,&x_1\geq1\text{ and }0\leq x_2<1,\\
  2/3,&x_2\geq1\text{ and }0\leq x_1<1,\\
  0,&\text{otherwise}
 \end{cases}
\]
\cite[p.~6]{DurrettPTEAuthor2019}. On the grid cut by the lines
\(x_i=0\) and \(x_i=1\), the value is \(0\) on every cell except the three
upper-right ones, which carry \(2/3\), \(1\) and \(2/3\).
\Cref{thm:p15-w02c-p15-002-negative-mass} shows that \(F\) satisfies
(i), (ii) and even (iii\(^*\)), while for \(0<\varepsilon\leq1\)
\[
 \Delta_{(1-\varepsilon,1]^2}F=1-\tfrac23-\tfrac23+0=-\tfrac13 .
\]
If a measure \(\mu\) had \(F\) as its distribution function, then by
\cref{thm:p15-w02c-p15-002-rectangle-inclusion-exclusion} the squares
\((1-\varepsilon,1]^2\), which decrease to the point \((1,1)\), would have
mass \(-1/3\); this is Durrett's computation
\(\mu(\{1,1\})=-1/3\), with \(\{1,1\}\) his notation for the point
\((1,1)\). The same limits give \(2/3\) at \((1,0)\) and at \((0,1)\),
so the three would-be point masses sum to \(1\) but one is negative. No
measure has this \(F\) as its distribution function.
\end{hbexample}

\begin{hblemma}{thm:p15-w02c-p15-002-negative-mass}{proved=annex:p15-w02c-p15-002-negative-mass}{The counterexample and a second one}
\begin{enumerate}
\item The function \(F\) of \cref{def:p15-w02c-p15-002-negative-mass}
      satisfies (i), (ii) and (iii\(^*\)); for \(0<\varepsilon\leq1\),
      \(\Delta_{(1-\varepsilon,1]^2}F=-1/3\) and
      \(\Delta_{(1-\varepsilon,1]\times(-\varepsilon,0]}F
       =\Delta_{(-\varepsilon,0]\times(1-\varepsilon,1]}F=2/3\). So
      (iv) fails, and no measure on \(\mathcal B(\mathbb R^2)\) has
      distribution function \(F\).
\item Let
      \(F_0(x_1,x_2)=\tfrac12\bigl(\mathbf1_{[0,\infty)}(x_1)
                     +\mathbf1_{[0,\infty)}(x_2)\bigr)\).
      Then \(F_0:\mathbb R^2\to[0,1]\) satisfies (i), (ii), (iii) and
      (iv), with \(\Delta_AF_0=0\) for every finite rectangle \(A\), but no
      probability measure \(\mu\) on \(\mathcal B(\mathbb R^2)\) satisfies
      \(\mu(A)=\Delta_AF_0\) for all finite rectangles. The function
      \(F_0\) fails (iii\(^*\)): \(F_0(t,0)=1/2\) for every \(t<0\).
\end{enumerate}
\end{hblemma}

Durrett's Theorem~1.1.11 reads: if \(F:\mathbb R^d\to[0,1]\) satisfies
(i)--(iv), then there is a unique probability measure \(\mu\) on
\((\mathbb R^d,\mathcal R^d)\) with \(\mu(A)=\Delta_AF\) for all finite
rectangles \cite[Theorem~1.1.11, p.~7]{DurrettPTEAuthor2019}. Part~(ii) of
\cref{thm:p15-w02c-p15-002-negative-mass} is a function satisfying
(i)--(iv) as printed for which no such probability measure exists, so the
theorem needs a hypothesis beyond the printed (iii). The proof in the
source first extends \(\mu\) from finite rectangles to rectangles with
infinite sides by monotonicity and then follows the one-dimensional
argument \cite[pp.~8--9]{DurrettPTEAuthor2019}; the extended values are
limits of \(\Delta_AF\) in which some vertices move to \(-\infty\) one
coordinate at a time, and those limits have the intended values only under
(iii\(^*\)). The handbook therefore states the theorem with (iii\(^*\)),
which by \cref{thm:p15-w02c-p15-002-rectangle-inclusion-exclusion} costs
nothing: every distribution function of a probability measure satisfies
it.

\begin{hbtheorem}{thm:p15-w02c-p15-002-rd-construction}{proved=annex:p15-w02c-p15-002-rd-construction}{Durrett's Theorem 1.1.11 with the strong tail condition}
\begin{enumerate}
\item Let \(F:\mathbb R^d\to\mathbb R\) satisfy (ii) and (iv). There is a
      unique measure \(\mu\) on \((\mathbb R^d,\mathcal B(\mathbb R^d))\)
      such that \(\mu(A)=\Delta_AF\) for every finite rectangle \(A\).
      It is finite on every bounded Borel set.
\item If in addition \(F\) takes values in \([0,1]\) and satisfies (i)
      and (iii\(^*\)), then \(\mu\) is a probability measure and
      \(F(x)=\mu((-\infty,x_1]\times\cdots\times(-\infty,x_d])\) for every
      \(x\in\mathbb R^d\).
\item Consequently a function \(F:\mathbb R^d\to[0,1]\) is the
      distribution function of a probability measure on
      \(\mathcal B(\mathbb R^d)\) if and only if it satisfies (i), (ii),
      (iii\(^*\)) and (iv), and the measure is then unique.
\end{enumerate}
\end{hbtheorem}

The proof is Durrett's in its two steps: additivity of \(\Delta_AF\) over
finite subdivisions, through subdivision into a grid, and countable
subadditivity, through right continuity and compactness of closed bounded
rectangles. It avoids rectangles with infinite sides by building the
measure on the boxes \((-n,n]^d\) and letting \(n\to\infty\); the tail
condition enters only in part~(ii).

\begin{hbexample}{def:p15-w02c-p15-002-product-distribution}{Durrett's Example 1.1.12: product form}
Let \(F(x)=\prod_{i=1}^dF_i(x_i)\), where each \(F_i:\mathbb R\to\mathbb R\)
is nondecreasing and right continuous
\cite[Example~1.1.12, p.~8]{DurrettPTEAuthor2019}. Then
\[
 \Delta_AF=\prod_{i=1}^d\bigl(F_i(b_i)-F_i(a_i)\bigr)\geq0
\]
for \(A=\prod_i(a_i,b_i]\)
(\cref{thm:p15-w02c-p15-002-product-difference}), and \(F\) is right
continuous, so part~(i) of \cref{thm:p15-w02c-p15-002-rd-construction}
gives a unique measure with these rectangle masses. When \(F_i(x)=x\) for
all \(i\) the masses are the volumes \(\prod_i(b_i-a_i)\); the product of
\(d\) copies of Lebesgue measure from
\cref{thm:p15-finite-product-measures} has the same values, so by the
uniqueness in part~(i) it is this measure, Lebesgue measure on
\(\mathbb R^d\), as the source states. When every \(F_i\) is the
distribution function of a real law, \(F\) satisfies (i) and (iii\(^*\))
as well, and the measure is the probability measure with distribution
function \(F\).
\end{hbexample}

\begin{hblemma}{thm:p15-w02c-p15-002-product-difference}{proved=annex:p15-w02c-p15-002-product-difference}{Differences of a product}
If \(F(x)=\prod_{i=1}^dF_i(x_i)\) with real functions \(F_i\), then
\(\Delta_AF=\prod_{i=1}^d(F_i(b_i)-F_i(a_i))\) for every finite rectangle
\(A=\prod_i(a_i,b_i]\). If each \(F_i\) is nondecreasing and right
continuous, \(F\) satisfies (ii) and (iv). If each \(F_i\) is moreover
the cumulative distribution function of a real random variable, \(F\)
maps into \([0,1]\) and satisfies (i) and (iii\(^*\)).
\end{hblemma}

\section{Membership notation for measurability, and indicators}
\label{sec:p15-w02c-p15-002-notation}

\begin{hbconvention}{def:p15-w02c-p15-002-membership-notation}{The notation \(X\in\mathcal F\)}
Let \((\Omega,\mathcal F)\) and \((S,\mathcal S)\) be measurable spaces and
\(X:\Omega\to S\). Durrett writes \(X\in\mathcal F\) to say that \(X\) is
\(\mathcal F\)-measurable, that is \(X^{-1}(B)\in\mathcal F\) for every
\(B\in\mathcal S\) \cite[p.~10]{DurrettPTEAuthor2019}. He uses the
notation again, for example writing \(G(q,\omega)\in\sigma(Y)\) for a
\(\sigma(Y)\)-measurable random variable \cite[p.~215]{DurrettPTEAuthor2019}.

The expression is not set membership: \(X\) is a function and
\(\mathcal F\) is a collection of subsets of \(\Omega\). Read as an
abbreviation, it means exactly
\[
 X\in\mathcal F\quad\Longleftrightarrow\quad\sigma(X)\subseteq\mathcal F ,
\]
by the definition of \(\sigma(X)\) in \cref{def:p15-measurable-map}. The
target sigma-algebra \(\mathcal S\) is implicit in the notation and must
be fixed by context; in the source it is \(\mathcal B(\mathbb R)\) for a
real function. For a set \(A\subseteq\Omega\) the expression
\(A\in\mathcal F\) keeps its ordinary meaning, and
\cref{thm:p15-w02c-p15-002-indicator-membership} shows that the two
readings agree on indicators. Elsewhere in this handbook measurability is
written out in words.
\index{measurable map!membership notation}
\end{hbconvention}

\begin{hblemma}{thm:p15-w02c-p15-002-indicator-membership}{proved=annex:p15-w02c-p15-002-indicator-membership}{Indicators and membership}
Let \(\mathcal F\) be a sigma-algebra on \(\Omega\) and
\(A\subseteq\Omega\), and let \(\mathbf1_A\) be one on \(A\) and zero on
\(A^c\), regarded as a map into \((\mathbb R,\mathcal B(\mathbb R))\).
Then
\[
 \sigma(\mathbf1_A)=\{\varnothing,A,A^c,\Omega\},
\]
so \(\mathbf1_A\) is \(\mathcal F\)-measurable exactly when
\(A\in\mathcal F\). In the membership notation,
\(\mathbf1_A\in\mathcal F\) if and only if \(A\in\mathcal F\).
\end{hblemma}

\begin{hbwarning}{def:p15-w02c-p15-002-indicator-terminology}{Indicator, not characteristic function}
The function \(\mathbf1_A\) of \cref{def:p15-simple-functions} is written
\(1_A\) by Durrett and is called the indicator function of \(A\). Durrett
notes that analysts call it the \emph{characteristic function} of \(A\),
while in probability that name belongs to a different object
\cite[p.~10]{DurrettPTEAuthor2019}: the function
\(t\mapsto\mathbb E e^{itX}\) of a random variable \(X\), defined in
\cref{def:p15-characteristic-function}. In this Part the phrase
\emph{characteristic function} always has the probabilistic meaning, and
\(\mathbf1_A\) is always called an indicator. Other Parts that use
\(\chi_A\) for the set-theoretic characteristic function do so in their
own declared scope.
\index{indicator of an event!terminology}\index{characteristic function!terminology}
\end{hbwarning}

\begin{hbmethod}{def:p15-w02c-p15-002-inverse-sampling}{Sampling through the generalized inverse}
\emph{Source form.} In the proof of his Theorem~1.2.2 Durrett takes
\(\Omega=(0,1)\) with the Borel sets and Lebesgue measure and puts
\(X(\omega)=\sup\{y:F(y)<\omega\}\); this is the map \(Q_F\) of
\cref{def:p15-quantile-map}. He calls it the \emph{inverse} of \(F\) and
writes \(F^{-1}\), although \(F\) need not be one-to-one or onto. He then
remarks that the scheme of this proof is a practical way to produce random
variables by computer: standard algorithms output a variable \(U\) with
the uniform distribution, and \(F^{-1}(U)\) then has distribution
function \(F\) \cite[Theorem~1.2.2, p.~11]{DurrettPTEAuthor2019}.

\emph{Inputs.} A function \(F\) as in \cref{def:p15-quantile-map}
(nondecreasing, right continuous, with limits \(0\) at \(-\infty\) and
\(1\) at \(\infty\)), and a random variable \(U\) on some probability
space with values in \((0,1)\) and \(P(U\leq u)=u\) for \(0<u<1\).

\emph{Procedure.} Output \(X=F^{-1}(U)=Q_F(U)\).

\emph{Conclusion.} \(X\) is a real random variable with
\(P(X\leq x)=F(x)\) for every real \(x\)
(\cref{thm:p15-w02c-p15-002-inverse-sampling}).

\emph{Side conditions.} The values \(0\) and \(1\) must be excluded from
the range of \(U\): at those points the defining supremum is \(-\infty\)
or can be \(+\infty\) (\cref{def:p15-cdf-existence-endpoints}). A
uniform variable that takes these values on a null set is first redefined
there, which changes no probability. Right continuity of \(F\) is needed
for the conclusion; the same warning shows a function without it whose
formal inverse produces the wrong distribution.

\emph{Worked example.} For the exponential distribution with rate
\(\lambda>0\), \(F(x)=0\) for \(x\leq0\) and \(F(x)=1-e^{-\lambda x}\) for
\(x\geq0\) \cite[Example~1.2.4, p.~12]{DurrettPTEAuthor2019}. For
\(0<u<1\) the set \(\{y:F(y)<u\}\) is the ray
\((-\infty,-\lambda^{-1}\log(1-u))\), so
\(F^{-1}(u)=-\lambda^{-1}\log(1-u)\) and
\(-\lambda^{-1}\log(1-U)\) is exponential with rate \(\lambda\). For a
distribution with a jump and a flat portion the inverse takes finitely
many values; \cref{def:p15-quantile-jump-flat} computes it.

\emph{Limitation.} The conclusion is a statement about an exactly uniform
\(U\). Whether the output of a particular generating algorithm may be
treated as such a variable is a different question, which this result does
not address.
\index{quantile!sampling}\index{distribution!simulation by inverse}
\end{hbmethod}

\begin{hbcorollary}{thm:p15-w02c-p15-002-inverse-sampling}{proved=annex:p15-w02c-p15-002-inverse-sampling}{The inverse applied to any uniform variable}
Let \(F\) be as in \cref{def:p15-quantile-map}, let
\((\Omega,\mathcal F,P)\) be a probability space, and let
\(U:\Omega\to(0,1)\) be a random variable with \(P(U\leq u)=u\) for
\(0<u<1\). Then \(X=Q_F(U)\) is a real random variable,
\[
 \{X\leq x\}=\{U\leq F(x)\}\quad\text{and}\quad P(X\leq x)=F(x)
 \qquad(x\in\mathbb R),
\]
and the law of \(X\) is the measure \(\mu_F\) of
\cref{thm:p15-cdf-existence}.
\end{hbcorollary}

\begin{hbconvention}{def:p15-w02c-p15-002-composition-notation}{The distribution as \(P\circ X^{-1}\)}
For a real random variable \(X\) on \((\Omega,\mathcal F,P)\), Durrett
calls the set function \(\mu(A)=P(X\in A)\), for Borel sets \(A\), the
\emph{distribution} of \(X\), and labels it \(\mu=P\circ X^{-1}\) in the
figure that accompanies the definition
\cite[sec.~1.2, Figure~1.3, p.~10]{DurrettPTEAuthor2019}. The composition is
one of set maps: \(X^{-1}\) sends a Borel set \(A\) to the event
\(X^{-1}(A)\in\mathcal F\), and \(P\) assigns that event its probability,
so
\[
 (P\circ X^{-1})(A)=P(X^{-1}(A))=P(X\in A)=\mathbb P_X(A)
 \qquad(A\in\mathcal B(\mathbb R)).
\]
Thus \(P\circ X^{-1}\) is the law \(\mathbb P_X\) of
\cref{def:p15-pushforward-law}, which that definition also calls the
distribution. Durrett checks countable additivity of \(\mu\) by pulling a
disjoint union back to a disjoint union of events
\cite[p.~10]{DurrettPTEAuthor2019}; the handbook's proof that
\(\mathbb P_X\) is a probability measure is \cref{thm:p15-pushforward-laws}.
The same notation applies to a measurable map into any measurable space
\((S,\mathcal S)\), with \(A\in\mathcal S\). The symbol \(X^{-1}\) here is
the preimage map on sets and exists for every function; it is not an
inverse function, and \(X\) need not be one-to-one.
\index{distribution!as P composed with X inverse}\index{pushforward measure!notation}
\end{hbconvention}

\begin{hbconvention}{def:p15-w02c-p15-002-equal-in-distribution}{Equal in distribution: \(X\stackrel{d}{=}Y\)}
Real random variables \(X\) and \(Y\), possibly defined on different
probability spaces, are \emph{equal in distribution} in Durrett's terms if
they induce the same distribution on \((\mathbb R,\mathcal B(\mathbb R))\)
\cite[p.~11]{DurrettPTEAuthor2019}. He writes
\[
 X\stackrel{d}{=}Y
\]
in displays, and \(X=_dY\) in running text, where the stacked symbol is
too tall \cite[p.~12]{DurrettPTEAuthor2019}. In the handbook's terms this
is the statement that \(X\) and \(Y\) have the \emph{same law}
(\cref{def:p15-pushforward-law}): \(\mathbb P_X=\mathbb P_Y\) on every
Borel set. Durrett notes that it holds exactly when the distribution
functions agree, \(P(X\leq x)=P(Y\leq x)\) for all \(x\)
\cite[pp.~11--12]{DurrettPTEAuthor2019}; this is the last statement of
\cref{thm:p15-cdf-laws}. The symbol \(X\stackrel{d}{=}Y\) has this meaning
in \cref{def:p15-weak-law-convergence}. The same words and symbols are used for random
elements of any common measurable value space, meaning equality of their
laws. Equality in distribution says nothing about the values at a given
outcome: if \(U\) is uniform on \((0,1)\), then \(U\stackrel{d}{=}1-U\),
since \(P(1-U\leq x)=P(U\geq1-x)=x=P(U\leq x)\) for \(0<x<1\) (both
distribution functions are \(0\) for \(x\leq0\) and \(1\) for
\(x\geq1\)), although \(U\neq1-U\) except at \(U=1/2\).
\index{equal in distribution}\index{law!equality of}
\end{hbconvention}

Durrett justifies the distribution-function criterion by his
Theorem~1.1.4 \cite[p.~11]{DurrettPTEAuthor2019},
which says that a Stieltjes measure function \(F\), that is a nondecreasing
and right-continuous real function, determines exactly one measure
\(\mu\) on the Borel sets of \(\mathbb R\) with \(\mu((a,b])=F(b)-F(a)\)
\cite[thm.~1.1.4, p.~3]{DurrettPTEAuthor2019}. The step from that theorem
to the criterion is not written in the source; the next lemma writes it,
with the handbook's version of Theorem~1.1.4,
\cref{thm:p15-stieltjes-construction}, in place of the source's.

\begin{hblemma}{thm:p15-w02c-p15-002-equal-in-distribution}{proved=annex:p15-w02c-p15-002-equal-in-distribution}{Equality in distribution through the uniqueness of Stieltjes measures}
Let \(X\) and \(Y\) be real random variables, on probability spaces
\((\Omega,\mathcal F,\mathbb P)\) and \((\Omega',\mathcal F',\mathbb P')\)
that may differ, with laws \(\mathbb P_X\), \(\mathbb P'_Y\) and
cumulative distribution functions \(F_X\), \(F_Y\).
\begin{enumerate}
\item \(F_X\) is a Stieltjes measure function in the sense of
      \cref{def:p15-stieltjes-intervals}, with endpoint values
      \(F_X(-\infty)=0\) and \(F_X(\infty)=1\), and
      \(\mathbb P_X(I(a,b))=F_X(b)-F_X(a)\) for all
      \(-\infty\leq a<b\leq\infty\). Hence \(\mathbb P_X=\mu_{F_X}\), the
      measure of \cref{thm:p15-stieltjes-construction}.
\item \(X\stackrel{d}{=}Y\) if and only if \(F_X(x)=F_Y(x)\) for every
      real \(x\).
\end{enumerate}
\end{hblemma}

The hypothesis that the laws are probability measures enters through
the endpoint values, and the criterion does not carry over to measures
of infinite mass as it stands. For Lebesgue measure \(\lambda\)
(\cref{thm:p15-lebesgue-unit-interval}) the ray \((-\infty,x]\) contains
\((x-n,x]\) for every \(n\), so \(\lambda((-\infty,x])\geq n\) for all
\(n\) and the function \(x\mapsto\lambda((-\infty,x])\) is identically
\(\infty\). The measure \(2\lambda\) has the same function, so the
function is not real-valued and does not determine the measure. Theorem~1.1.4 is then used with \(F(x)=x\), a function
chosen to match the interval lengths, not one read off from rays.

\section{A discrete law whose jumps are dense}
\label{sec:p15-w02c-p15-002-dense-jumps}

A discrete law in Durrett's sense gives probability one to a countable set
(\cref{def:p15-w02c-p15-014-point-mass-law}). After the point mass at~\(0\),
his simplest discrete law, he gives an example to show that the
distribution function of a discrete law need not look like a staircase
\cite[p.~13, Example~1.2.9, p.~14]{DurrettPTEAuthor2019}.

\begin{hbexample}{def:p15-w02c-p15-002-dense-jumps}{Durrett's Example 1.2.9: dense discontinuities}
Let \(q_1,q_2,\ldots\) be an enumeration of the rational numbers, that is,
a sequence in which every rational number occurs exactly once. Let
\(\alpha_i>0\) with \(\sum_{i=1}^\infty\alpha_i=1\), and put
\[
 F(x)=\sum_{i=1}^\infty\alpha_i\mathbf1_{[q_i,\infty)}(x)
 \qquad(x\in\mathbb R),
\]
where \(\mathbf1_{[\theta,\infty)}(x)\) is \(1\) for \(x\geq\theta\) and
\(0\) otherwise \cite[Example~1.2.9, p.~14]{DurrettPTEAuthor2019}. (The
normalization is printed as \(\sum_{i=1}^\infty\alpha_1=1\), which cannot
hold for a positive constant term; the sum of the \(\alpha_i\) is meant.)
For instance \(\alpha_i=2^{-i}\) satisfies the conditions.

\Cref{thm:p15-w02c-p15-002-dense-jumps} shows what the example exhibits:
\(F\) is the distribution function of the discrete law
\(\mu=\sum_i\alpha_i\delta_{q_i}\), which gives probability one to
\(\mathbb Q\); \(F\) jumps by \(\alpha_i\) at \(q_i\), so every nonempty
open interval contains infinitely many jumps; \(F\) is continuous at every
irrational number; and \(F\) is strictly increasing, so it is constant on
no interval. Hence \(F\) has no interval on which it is a step function,
although the law is discrete. Reading the law off the graph therefore
requires all of the countably many jumps, not finitely many of them.

The enumeration matters only through its range. If a rational number were
allowed to occur several times, the jump of \(F\) there would be the sum of
the \(\alpha_i\) over its occurrences, and the conclusions would be
unchanged.
\index{distribution function!dense discontinuities}\index{distribution!discrete}
\end{hbexample}

\begin{hblemma}{thm:p15-w02c-p15-002-dense-jumps}{proved=annex:p15-w02c-p15-002-dense-jumps}{The law behind Example 1.2.9}
Let \(q_i\), \(\alpha_i\) and \(F\) be as in
\cref{def:p15-w02c-p15-002-dense-jumps}, and for Borel \(B\subseteq\mathbb R\)
put
\[
 \mu(B)\defeq\sum_{i:\,q_i\in B}\alpha_i .
\]
\begin{enumerate}
\item \(\mu\) is a probability measure on \((\mathbb R,\mathcal B(\mathbb R))\)
      with \(\mu(\mathbb R\setminus\mathbb Q)=0\), so it is discrete with
      countable carrier \(\mathbb Q\); and
      \(F(x)=\mu((-\infty,x])\) for every real \(x\). In particular \(F\) is
      nondecreasing and right continuous, with limits \(0\) at \(-\infty\)
      and \(1\) at \(\infty\).
\item \(F(q_i)-F(q_i-)=\alpha_i\) for every \(i\), and
      \(F(x)-F(x-)=0\) for every irrational \(x\). The set of
      discontinuities of \(F\) is exactly \(\mathbb Q\), and every nonempty
      open interval contains infinitely many of them.
\item \(F\) is strictly increasing: \(x<y\) implies \(F(x)<F(y)\).
\end{enumerate}
\end{hblemma}

\section{The extended line, the convergence set and the limit convention}
\label{sec:p15-w02c-p15-002-extended}

\Cref{def:p15-extended-random-variables} uses \(\overline{\mathbb R}\)
with the sigma-algebra generated by the rays \([-\infty,a]\), works with
functions defined on all of \(\Omega\), and defines almost sure
convergence to a specified real limit. The source differs on each point
\cite[pp.~16--17]{DurrettPTEAuthor2019}. Its forms are as follows.

\begin{hbdefinition}{def:p15-w02c-p15-002-extended-line}{Durrett's extended line, random variables on an event, and almost sure convergence}
\begin{enumerate}
\item \emph{The extended line.} \(\mathbf R^*\defeq[-\infty,\infty]\),
      the real line with a least element \(-\infty\) and a greatest
      element \(\infty\) adjoined; as an ordered set it is
      \(\overline{\mathbb R}\). It carries the topology generated by the
      intervals \([-\infty,a)\), \((a,b)\) and \((b,\infty]\) for real
      \(a,b\), and \(\mathcal R^*\) is the Borel sigma-algebra of that
      topology, the sigma-algebra generated by its open sets. The pair
      \((\mathbf R^*,\mathcal R^*)\) is the \emph{extended real line}
      \cite[p.~17]{DurrettPTEAuthor2019}.
\item \emph{Random variables on an event.} For \(D\in\mathcal F\), a
      function \(X:D\to\mathbf R^*\) is a \emph{random variable} in the
      source's extended sense if \(X^{-1}(B)=\{\omega\in D:X(\omega)\in
      B\}\in\mathcal F\) for every \(B\in\mathcal R^*\)
      \cite[p.~17]{DurrettPTEAuthor2019}. The domain \(D\) is part of the
      data. By contrast, every measurable map in
      \cref{def:p15-measurable-map} is defined on all of \(\Omega\).
\item \emph{The convergence set.} For real random variables
      \(X_1,X_2,\ldots\), Durrett defines
      \(\Omega_o\) as the set of outcomes \(\omega\) at which
      \(\lim_{n\to\infty}X_n(\omega)\) exists
      \cite[p.~16]{DurrettPTEAuthor2019}. Here the limit is read as a
      real number:
      \[
       \Omega_o\defeq\{\omega:\lim_{n\to\infty}X_n(\omega)
                        \text{ exists in }\mathbb R\}.
      \]
      The larger set on which the limit exists in \(\mathbf R^*\) is
      written out in full, as
      \(\{\liminf_nX_n=\limsup_nX_n\}\), when it is meant.
\item \emph{Almost sure convergence without a prescribed limit.} The
      sequence \(X_n\) \emph{converges almost surely} in the source's
      sense if \(\mathbb P(\Omega_o)=1\)
      \cite[pp.~16--17]{DurrettPTEAuthor2019}. No limit variable is named
      in this definition. Measure theory calls the same notion
      convergence almost everywhere.
\item \emph{The limit convention.} To have a limit defined on the whole
      space, Durrett sets
      \[
       X_\infty\defeq\limsup_{n\to\infty}X_n ,
      \]
      a function on all of \(\Omega\) that may take the values
      \(\pm\infty\) \cite[p.~17]{DurrettPTEAuthor2019}.
\end{enumerate}
The source describes \(\Omega_o\) also as the set where
\(\limsup_nX_n-\liminf_nX_n=0\). Where both are infinite with the same
sign this difference is the undefined expression \(\infty-\infty\), so
the handbook uses the descriptions in
\cref{thm:p15-w02c-p15-002-convergence-set} instead; on the set where
\(\limsup_nX_n\) and \(\liminf_nX_n\) are both finite, the source's
description is exact.
\index{extended real line}\index{random variable!on an event}
\index{convergence!almost surely!without prescribed limit}
\end{hbdefinition}

\begin{hbtheorem}{thm:p15-w02c-p15-002-order-borel}{proved=annex:p15-w02c-p15-002-order-borel}{Borel sets of the extended line}
The sigma-algebra \(\mathcal R^*\) equals
\(\sigma(\{[-\infty,a]:a\in\mathbb R\})\), the sigma-algebra placed on
\(\overline{\mathbb R}\) in \cref{def:p15-extended-random-variables}. More
precisely, every open set of \(\mathbf R^*\) is a countable union of
intervals \([-\infty,q)\), \((q,r)\) and \((r,\infty]\) with rational
\(q,r\). Hence a function \(Z:\Omega\to\mathbf R^*\) is
\(\mathcal F/\mathcal R^*\)-measurable exactly when it is an extended real
random variable in the sense of
\cref{def:p15-extended-random-variables}, and every result proved there
for \(\overline{\mathbb R}\) holds verbatim for \((\mathbf R^*,\mathcal R^*)\).
\end{hbtheorem}

\begin{hblemma}{thm:p15-w02c-p15-002-domain-extension}{proved=annex:p15-w02c-p15-002-domain-extension}{Random variables on an event}
Let \(D\in\mathcal F\), \(X:D\to\mathbf R^*\), and \(c\in\mathbf R^*\).
Define \(X^c:\Omega\to\mathbf R^*\) by \(X^c=X\) on \(D\) and \(X^c=c\) on
\(\Omega\setminus D\). The following are equivalent.
\begin{enumerate}
\item \(X\) is a random variable on \(D\) in the sense of
      \cref{def:p15-w02c-p15-002-extended-line}.
\item \(X\) is measurable from \((D,\mathcal F_D)\) to
      \((\mathbf R^*,\mathcal R^*)\), where
      \(\mathcal F_D=\{A\in\mathcal F:A\subseteq D\}\).
\item \(X^c\) is an extended real random variable on \(\Omega\).
\end{enumerate}
The collection \(\mathcal F_D\) is a sigma-algebra on \(D\), equal to
\(\{D\cap A:A\in\mathcal F\}\).
\end{hblemma}

The lemma converts every statement about random variables on an event
into one about functions on all of \(\Omega\), which is the form used in
the rest of the handbook; the choice of the value \(c\) off \(D\) does not
affect measurability.

\begin{hbtheorem}{thm:p15-w02c-p15-002-convergence-set}{proved=annex:p15-w02c-p15-002-convergence-set}{The convergence set is an event, and the two forms of almost sure convergence}
Let \(X_1,X_2,\ldots\) be real random variables, and write
\(L=\liminf_nX_n\) and \(U=\limsup_nX_n=X_\infty\).
\begin{enumerate}
\item \(L\leq U\) everywhere. A real sequence has a limit in
      \(\mathbf R^*\) exactly when its lower and upper limits are equal,
      and the limit is then their common value.
\item The sets
      \[
       \{L<U\}=\bigcup_{q,r\in\mathbb Q,\ q<r}\bigl(\{L<q\}\cap\{U>r\}\bigr),
      \]
      \[
       \{L=U\}=\Omega\setminus\{L<U\},\qquad
       \Omega_o=\{L=U\}\cap\bigcup_{m\geq1}\{-m<U<m\}
      \]
      are events. On \(\{L=U\}\) the limit of \(X_n\) in \(\mathbf R^*\)
      exists and equals \(X_\infty\); on \(\Omega_o\) it is real.
\item \(X_\infty\) is an extended real random variable, and the
      restriction of \(X_\infty\) to \(D=\{L=U\}\) is a random variable on
      \(D\) in the sense of \cref{def:p15-w02c-p15-002-extended-line}.
\item If \(\mathbb P(\Omega_o)=1\), let \(Y=X_\infty\) on \(\Omega_o\) and
      \(Y=0\) off \(\Omega_o\). Then \(Y\) is a real random variable and
      \(X_n\) converges almost surely to \(Y\) in the sense of
      \cref{def:p15-extended-random-variables}.
\item Conversely, if \(X_n\) converges almost surely to a real random
      variable \(X\) in the sense of
      \cref{def:p15-extended-random-variables}, then
      \(\mathbb P(\Omega_o)=1\) and \(X=X_\infty\) on the probability-one
      event \(\{X_n\to X\}\). Any two such limits agree outside a null
      event.
\end{enumerate}
\end{hbtheorem}

The source's definition and the handbook's therefore describe the same
sequences, and the limit is determined up to a null event; the source's
\(X_\infty\) is one representative, defined everywhere but possibly
infinite off \(\Omega_o\). For the partial sums \(S_n\) of a random series,
\(\Omega_o\) is the convergence event \(C_X\) of
\cref{thm:p15-series-event-calculus}, where it is written through the
Cauchy criterion.

\begin{hbexample}{def:p15-w02c-p15-002-convergence-examples}{Three convergence sets}
Work on the uniform probability space
\((\Omega_U,\mathcal F_U,\mathbb U)\) on \((0,1)\) of
\cref{thm:p15-lebesgue-unit-interval}.
\begin{enumerate}
\item \(X_n(\omega)=n\mathbf1_{(0,1/n)}(\omega)\). For each \(\omega\),
      \(X_n(\omega)=0\) as soon as \(1/n\leq\omega\), so the limit is \(0\)
      at every outcome. Here \(\Omega_o=\Omega_U\), the sequence converges
      almost surely, and \(X_\infty=0\), although \(\mathbb E X_n=1\) for
      every \(n\).
\item \(X_n(\omega)=n\omega\). For every \(\omega\in(0,1)\) the sequence
      tends to \(\infty\). Hence \(\{L=U\}=\Omega_U\) but
      \(\Omega_o=\varnothing\): the limit exists in \(\mathbf R^*\) at
      every outcome, yet the sequence does not converge almost surely, and
      \(X_\infty=\infty\) everywhere. This is the case in which the
      source's difference \(U-L\) is \(\infty-\infty\).
\item \(X_n(\omega)=(-1)^n\). Then \(L=-1<1=U\) everywhere, and
      \(\Omega_o=\varnothing\); nevertheless \(X_\infty=1\) is defined
      everywhere. The convention assigns a value even where no limit
      exists, which is why it is used only on \(\Omega_o\) or with an
      explicit statement that it is a \(\limsup\).
\end{enumerate}
\end{hbexample}

\section{Uniqueness for measures that are finite on an exhausting sequence}
\label{sec:p15-w02c-p15-002-sigma-finite-uniqueness}

\Cref{thm:p15-measure-uniqueness} proves uniqueness for two probability
measures that agree on a generating pi-system, and
\cref{thm:p15-premeasure-extension} proves uniqueness of the extension of a
sigma-finite premeasure from an algebra. Durrett's Theorem~A.1.5 is the
general pi-system form for measures that need not be finite
\cite[Theorem~A.1.5, pp.~457--458]{DurrettPTEAuthor2019}.

\begin{hbtheorem}{thm:p15-w02c-p15-002-sigma-finite-uniqueness}{proved=annex:p15-w02c-p15-002-sigma-finite-uniqueness}{Durrett's Theorem A.1.5: finite exhaustion in a pi-system}
Let \(\mathcal P\) be a pi-system on \(\Omega\). Let \(\nu_1\) and \(\nu_2\)
be measures on sigma-algebras \(\mathcal F_1\) and \(\mathcal F_2\) on
\(\Omega\), each containing \(\sigma(\mathcal P)\). Suppose that
\(\nu_1(A)=\nu_2(A)\) for every \(A\in\mathcal P\), and that there are sets
\(A_n\in\mathcal P\) with \(A_n\subseteq A_{n+1}\),
\(\bigcup_nA_n=\Omega\) and \(\nu_i(A_n)<\infty\). Then
\(\nu_1(B)=\nu_2(B)\) for every \(B\in\sigma(\mathcal P)\).
\end{hbtheorem}

Durrett derives Theorem~A.1.5 from his Theorem~A.1.4, which he calls
Dynkin's \(\pi\)-\(\lambda\) theorem: a lambda-system that contains a
pi-system \(\mathcal P\) contains \(\sigma(\mathcal P)\). He names it the key
to his proof that the extension of a measure is unique, and says it will be
used many times in the book \cite[Theorem~A.1.4, p.~456]{DurrettPTEAuthor2019}.
His pi-systems and lambda-systems, defined on the same page, are those of
\cref{def:p15-pi-lambda}: a lambda-system contains \(\Omega\), is closed
under \(B\setminus A\) for members \(A\subseteq B\), and is closed under
increasing unions. His Theorem~A.1.4 is the first statement of
\cref{thm:p15-measure-uniqueness}, and the proof of
\cref{thm:p15-w02c-p15-002-sigma-finite-uniqueness} applies it in that form.

The exhausting sequence must lie in \(\mathcal P\). For probability
measures \cref{thm:p15-measure-uniqueness} needs no such sequence, because
the total masses are both one; if \(\Omega\in\mathcal P\) the probability
case is the special case \(A_n=\Omega\).

\begin{hbexample}{def:p15-w02c-p15-002-uniqueness-failure}{Agreement on a generating pi-system without finite exhaustion}
Let \(\Omega=\mathbb Q\) with all subsets measurable, and let \(\nu_1\) and
\(\nu_2\) be the countable atomic measures of
\cref{thm:p15-countable-atomic} with weights \(1\) and \(2\) at every
point. Let
\(\mathcal P=\{\varnothing\}\cup\{\mathbb Q\cap(a,b]:a<b\text{ real}\}\).
The intersection of two members is empty or again of this form, so
\(\mathcal P\) is a pi-system. Each nonempty member contains infinitely
many rationals, because between any two reals there is a rational, so
both measures give it the value \(\infty\); they also agree on
\(\varnothing\). The sets \(\mathbb Q\cap(-n,n]\) belong to \(\mathcal P\)
and increase to \(\Omega\), but have infinite measure. Since
\(\{0\}=\bigcap_n\mathbb Q\cap(-1/n,0]\in\sigma(\mathcal P)\) and
\(\nu_1(\{0\})=1\neq2=\nu_2(\{0\})\), the measures differ on
\(\sigma(\mathcal P)\). The finiteness condition on the exhausting sequence
cannot be dropped.
\end{hbexample}

\section{Nice spaces}
\label{sec:p15-w02c-p15-002-nice}

Durrett's definition of a \emph{nice} measurable space, which asks for a
one-to-one map \(\varphi\) into \(\mathbb R\) with \(\varphi\) and
\(\varphi^{-1}\) measurable and has no clause requiring \(\varphi(S)\) to
be a Borel set, is stated in \cref{def:p15-w02c-p15-004-nice}, with his
remark that such spaces are often called standard Borel spaces
\cite[p.~53]{DurrettPTEAuthor2019}. The same section gives the exact
relation to Borel real codings (\cref{thm:p15-w02c-p15-004-nice-coding}),
his Theorem~2.1.22 in his terms
(\cref{thm:p15-w02c-p15-004-durrett-2122}), his statement that the
extension theorem fails for arbitrary measurable spaces with its
attributions (\cref{def:p15-w02c-p15-004-extension-failure}), the status of
his claim that the generalization of Kolmogorov's theorem to nice spaces
is trivial (\cref{thm:p15-w02c-p15-004-nice-extension}), and his remark on
the cardinality of Borel sets \cite[p.~54]{DurrettPTEAuthor2019}, proved
in \cref{sec:p15-w02c-p15-004-borel-cardinality}. This chapter uses that
definition, and a map with its properties is called a \emph{nice map}.

One fact is needed below. By part~(iii) of
\cref{thm:p15-w02c-p15-004-nice-coding}, a nice map \(\varphi\) carries
every \(A\in\mathcal S\) to a Borel subset of \(\mathbb R\) exactly when
\(\varphi(S)\) is Borel. In his existence proof for regular conditional
distributions on a nice space, Durrett passes from \(A\in\mathcal S\) to
the set \(\varphi(A)=(\varphi^{-1})^{-1}(A)\) and treats it as a Borel
set \cite[Theorem~4.1.17, p.~215]{DurrettPTEAuthor2019}. That step is
therefore valid for every \(A\) exactly when the image of \(\varphi\) is
Borel. \Cref{thm:p15-standard-conditional-law} and
\cref{thm:p15-coded-space-extension} are proved for Borel real codings;
\cref{sec:p15-w02c-p15-002-nice-rcd} settles the source's
Theorems~4.1.17 and~4.1.18 under the definition as printed.

\section{Borel's paradox}
\label{sec:p15-w02c-p15-002-borel-paradox}

\begin{hbexample}{def:p15-w02c-p15-002-borel-paradox}{Borel's paradox in the source's framing}
Durrett gives the sphere example the name \emph{Borel's paradox}
\cite[Example~4.1.8, p.~210]{DurrettPTEAuthor2019}. A point is chosen at
random on the earth, with longitude \(\theta\) and latitude \(\varphi\).
The customary ranges are \(\theta\in[0,2\pi)\) and
\(\varphi\in(-\pi/2,\pi/2]\). The source points out that one may equally
take \(\theta\in[0,\pi)\) and \(\varphi\in(-\pi,\pi]\): the new longitude
then selects the great circle through the poles on which the point lies,
and \(\varphi\) is the angle along that circle.

\emph{The apparent paradox.} Since the point is uniform on the sphere and
the great circle through the poles is chosen uniformly, one might expect
the angle along the circle to be uniform as well. In the source's account,
\(\theta\) is uniform but the angle is not.

\emph{The source's resolution.} In both coordinate systems the angle is
independent of the longitude, so its conditional distribution given the
longitude is its unconditional distribution. That distribution is not
uniform, because more of the sphere's surface lies near the equator than
near the poles \cite[p.~210]{DurrettPTEAuthor2019}.

In the handbook's notation these coordinates are \(\Theta,\Phi\) and
\(\Theta',\Psi\) of \cref{def:p15-sphere-coordinate-law}, and
\cref{thm:p15-conditioning-coordinate-examples} proves the independence
and the densities. \Cref{thm:p15-w02c-p15-002-borel-paradox} turns the
source's resolution into statements about regular conditional laws and
measures the concentration near the equator.
\index{Borel's paradox}
\end{hbexample}

\begin{hbcorollary}{thm:p15-w02c-p15-002-borel-paradox}{proved=annex:p15-w02c-p15-002-borel-paradox}{The conditional law of the angle along the circle}
In the spherical model of \cref{def:p15-sphere-coordinate-law}:
\begin{enumerate}
\item The constant kernel \(K(\omega,A)=\mathbb P(\Psi\in A)\),
      \(A\in\mathcal B(\mathbb R)\), is a regular conditional law of
      \(\Psi\) given \(\sigma(\Theta')\), and the constant kernel
      \(\mathbb P(\Phi\in\cdot)\) is one of \(\Phi\) given
      \(\sigma(\Theta)\).
\item Every regular conditional law \(K'\) of \(\Psi\) given
      \(\sigma(\Theta')\) satisfies \(K'(\omega,\cdot)=\mathbb P(\Psi\in\cdot)\)
      for all \(\omega\) outside one null event. In particular, outside a
      null event \(K'(\omega,\cdot)\) is not the uniform law on
      \((-\pi,\pi]\).
\item \(\mathbb P(|\Phi|\leq\pi/6)=1/2\) and
      \(\mathbb P(|\Phi|\geq\pi/3)=1-\sqrt3/2\), while the corresponding
      fractions of the latitude range \((-\pi/2,\pi/2]\) are \(1/3\) and
      \(1/3\). Likewise \(\mathbb P(|\Psi|\leq\pi/6)=1/4\), while
      \((-\pi/6,\pi/6]\) is the fraction \(1/6\) of \((-\pi,\pi]\).
\end{enumerate}
\end{hbcorollary}

Thus half of the sphere's area lies within latitude \(\pi/6\) of the
equator, a band covering one third of the latitude range, and the polar
caps beyond latitude \(\pi/3\), also one third of the range, carry less
than \(14\) per cent of the area. The density \(|\cos\psi|/4\) of
\(\Psi\) is largest at \(\psi=0\) and \(\psi=\pi\), where the great circle
crosses the equator, and vanishes at \(\psi=\pm\pi/2\), the poles. No
conditional probability given the null event of a single longitude has
been used: the regular conditional law is a kernel, determined up to a
null set of longitudes, and the second part of
\cref{thm:p15-conditioning-coordinate-examples} shows how a different
choice of conditioning variable changes the law attached to the same null
event.

\section{Why regular conditional distributions can fail to exist}
\label{sec:p15-w02c-p15-002-rcd-existence}

For each fixed measurable set the scalar conditional probability exists
(\cref{thm:p15-conditional-existence}), and the countable additivity of
those versions holds outside a null set that depends on the sets
involved.

\begin{hblemma}{thm:p15-w02c-p15-002-ae-additivity}{proved=annex:p15-w02c-p15-002-ae-additivity}{Countable additivity outside a null set}
Let \(X:(\Omega,\mathcal F_0,P)\to(S,\mathcal S)\) be measurable and
\(\mathcal G\subseteq\mathcal F_0\) a sub-sigma-algebra. For each
\(A\in\mathcal S\) choose a version \(g_A\) of \(P(X\in A\mid\mathcal G)\).
\begin{enumerate}
\item For every sequence of pairwise disjoint \(A_1,A_2,\ldots\in\mathcal S\)
      there is a \(\mathcal G\)-measurable null set \(N\) such that
      \[
       g_{\bigcup_nA_n}(\omega)=\sum_{n\geq1}g_{A_n}(\omega)
                                    \qquad(\omega\notin N).
      \]
\item For any countable family of such sequences a single
      \(\mathcal G\)-measurable null set serves for all of them.
\end{enumerate}
\end{hblemma}

\begin{hbwarning}{def:p15-w02c-p15-002-rcd-nonexistence}{Exceptional sets that accumulate}
Durrett records that regular conditional distributions do not always exist,
that the first example is due to Dieudonn\'e (1948), and that Doob (1953),
p.~624, and Faden (1985) treat later developments
\cite[p.~214]{DurrettPTEAuthor2019}. Those examples are reported only as
the source names them. A probability space with no regular conditional
distribution, constructed from the measure theory of this Part and not
offered as any of them, is given in \cref{thm:p15-w02c-p15-008-no-rcp}.

The source's explanation of the difficulty is the one isolated in
\cref{thm:p15-w02c-p15-002-ae-additivity}. For disjoint
\(A_1,A_2,\ldots\), linearity and monotone convergence of conditional
expectation give
\(P(X\in\bigcup_nA_n\mid\mathcal G)=\sum_nP(X\in A_n\mid\mathcal G)\)
almost surely \cite[p.~214]{DurrettPTEAuthor2019}; in the handbook these
are \cref{thm:p15-conditional-basic} and part~(i) of
\cref{thm:p15-conditional-limits}. Each such identity has its own
exceptional null set. When \(\mathcal S\) contains uncountably many
countable disjoint families, the union of the exceptional sets need not
be null, and then no choice of versions need give a probability measure
\(A\mapsto g_A(\omega)\) for almost every \(\omega\). Countable unions of
null sets are null (\cref{thm:p15-measure-laws}); uncountable unions need
not be.

The existence proofs in the handbook avoid the accumulation by
determining the whole kernel from countably many versions: for a real
random variable, the versions of \(P(W\leq q\mid\mathcal G)\) at rational
\(q\) (\cref{thm:p15-real-conditional-law}), and for a space with a Borel
real coding, the same construction transported through the coding
(\cref{thm:p15-standard-conditional-law}). Only countably many exceptional
sets are discarded there.
\index{regular conditional distribution!nonexistence}
\end{hbwarning}

\section{Regular conditional distributions on nice spaces}
\label{sec:p15-w02c-p15-002-nice-rcd}

Durrett's definition of a regular conditional distribution is as follows
\cite[p.~214]{DurrettPTEAuthor2019}. Let \((\Omega,\mathcal F,P)\) be a
probability space, \(X:(\Omega,\mathcal F)\to(S,\mathcal S)\) measurable
and \(\mathcal G\subseteq\mathcal F\) a sigma-field. A function
\(\mu:\Omega\times\mathcal S\to[0,1]\) is a \emph{regular conditional
distribution} for \(X\) given \(\mathcal G\) if (i) for each \(A\),
\(\omega\mapsto\mu(\omega,A)\) is a version of \(P(X\in A\mid\mathcal G)\),
and (ii) for almost every \(\omega\), \(A\mapsto\mu(\omega,A)\) is a
probability measure on \((S,\mathcal S)\). A regular conditional law in
the sense of \cref{def:p15-regular-conditional-law} is one in which (ii)
holds for every \(\omega\); it is in particular a regular conditional
distribution in the source's sense.

The source's two existence theorems assume a nice space in the sense of
\cref{def:p15-w02c-p15-004-nice}:
\begin{itemize}
\item \emph{Theorem~4.1.17}: regular conditional distributions exist if
      \((S,\mathcal S)\) is nice \cite[p.~215]{DurrettPTEAuthor2019};
\item \emph{Theorem~4.1.18}: if \(X\) and \(Y\) take values in a nice space
      \((S,\mathcal S)\) and \(\mathcal G=\sigma(Y)\), there is a function
      \(\mu:S\times\mathcal S\to[0,1]\) such that (i) for each \(A\),
      \(\mu(Y(\omega),A)\) is a version of \(P(X\in A\mid\mathcal G)\), and
      (ii) for almost every \(\omega\), \(A\mapsto\mu(Y(\omega),A)\) is a
      probability measure on \((S,\mathcal S)\)
      \cite[p.~215]{DurrettPTEAuthor2019}.
\end{itemize}
Both proofs apply the real construction to \(\varphi(X)\) and then set
\(\mu(\omega,A)=\nu(\omega,\varphi(A))\), using that
\(\varphi(A)=(\varphi^{-1})^{-1}(A)\) is a Borel set. By part~(iii) of
\cref{thm:p15-w02c-p15-004-nice-coding} that holds for every
\(A\in\mathcal S\) exactly when \(\varphi(S)\) is Borel, which is the case
proved in \cref{thm:p15-standard-conditional-law,thm:p15-value-conditional-law}.
The next two results settle the source's theorems for the definition of a
nice space as printed.

\begin{hbtheorem}{thm:p15-w02c-p15-002-nice-rcd}{proved=annex:p15-w02c-p15-002-nice-rcd}{Theorems 4.1.17 and 4.1.18 for nice spaces whose law is carried by a Borel set}
Let \((S,\mathcal S)\) be nice with nice map \(\varphi\), let
\(X:(\Omega,\mathcal F,P)\to(S,\mathcal S)\) be measurable, and suppose
there is a Borel set \(E\subseteq\varphi(S)\) with
\(P(\varphi(X)\in E)=1\). This holds with \(E=\varphi(S)\) whenever
\(\varphi(S)\) is Borel.
\begin{enumerate}
\item For every sub-sigma-algebra \(\mathcal G\subseteq\mathcal F\), \(X\)
      has a regular conditional law given \(\mathcal G\), hence a regular
      conditional distribution in the source's sense.
\item For every measurable \(Y:\Omega\to(T,\mathcal T)\) there is a
      probability kernel \(K\) from \((T,\mathcal T)\) to
      \((S,\mathcal S)\) such that \(K(Y,A)\) is a version of
      \(P(X\in A\mid\sigma(Y))\) for each \(A\in\mathcal S\). With
      \((T,\mathcal T)=(S,\mathcal S)\) and \(\mu=K\), \(\mu\) satisfies
      (i) and (ii) of Theorem~4.1.18, with (ii) for every \(\omega\).
\end{enumerate}
\end{hbtheorem}

\begin{hbtheorem}{thm:p15-w02c-p15-002-nice-rcd-failure}{proved=annex:p15-w02c-p15-002-nice-rcd-failure}{Without a Borel carrier both theorems can fail}
Assume the axiom of choice. There are a nice space \((S,\mathcal S)\), a
probability space \((\Omega,\mathcal F,P)\) and measurable
\(X,Y:\Omega\to S\) such that \(X\) has no regular conditional
distribution given \(\sigma(Y)\) in the source's sense. In particular no
function \(\mu:S\times\mathcal S\to[0,1]\) has properties (i) and (ii) of
Theorem~4.1.18 for this pair. For every nice map \(\varphi\) of this space,
no Borel subset of \(\varphi(S)\) carries the law of \(\varphi(X)\); in
particular \(\varphi(S)\) is not Borel.
\end{hbtheorem}

The example uses the probability space \((\Omega_U,\mathcal F_H,P_H)\) of
\cref{def:p15-w02c-p15-008-model}. It is built from a set
\(H_U\subseteq(0,1)\) such that every Borel subset of \(H_U\) and every
Borel subset of \((0,1)\setminus H_U\) has length zero
(\cref{thm:p15-w02c-p15-008-thick-set}); \(P_H\) extends the uniform
probability \(\mathbb U\) and gives \(B\cap H_U\) probability
\(\mathbb U(B)/2\) for Borel \(B\subseteq(0,1)\)
(\cref{thm:p15-w02c-p15-008-model}). Given the Borel sets, \(H_U\) then
has conditional probability \(1/2\), while a regular conditional
distribution would have to put point mass at the outcome itself, and no
regular conditional distribution of the identity map exists
(\cref{thm:p15-w02c-p15-008-no-rcp}). The space \(S\) below consists of two
copies of \((0,1)\), one carrying \(\mathcal F_H\) and observed by \(X\),
the other carrying the Borel sets and observed by \(Y\).

The positive theorem and the example together give the exact status of the
source's two theorems under its printed definition: they hold whenever
the law of \(\varphi(X)\) is carried by a Borel subset of the image, in
particular for every nice space with Borel image, and they can fail
otherwise. Every space that the source shows to be nice, the Borel subsets
of complete separable metric spaces, has a nice map with Borel image by
part~(ii) of \cref{thm:p15-w02c-p15-004-durrett-2122}, so the source's applications are
covered.
